%% kapa: exam=1 round=34 year=2023 period=1 form=A subject=economics subject_ko=경제학원론 seq=18 number=58
%% source: Q-Net 감정평가사 기출문제 (cst00309) artlSeq=5211789 file=raw/1008/5211789/2023년도 제34회 감정평가사 1차시험 1교시.pdf
%% conversion: kordoc 4.24.51 -> tools/gen1.py (auto); review=proofread-sonnet; flags=none
\begin{question}{58}
정보재(information goods)의 기본적인 특성에 관한 설명으로 옳은 것을 모두 고른 것은?
\begin{framedtext}
ㄱ. 상품에 포함된 정보가 상품으로서의 특성을 결정하는 것을 정보재라 한다. \par
ㄴ. 정보재는 초기 개발비용이 크고 한계비용이 0에 가깝기 때문에 규모의 불경제가 일어난다. \par
ㄷ. 정보재에는 쏠림현상(tipping)과 같은 네트워크효과가 나타난다. \par
ㄹ. 정보재의 경우 무료견본을 나눠주는 것은 잠김효과(lock-in effect)를 노린 마케팅 전략이다. \par
\end{framedtext}
\begin{choices}
  \item ㄱ, ㄹ
  \item ㄴ, ㄷ
  \item ㄱ, ㄴ, ㄷ
  \item ㄱ, ㄷ, ㄹ
  \item ㄴ, ㄷ, ㄹ
\end{choices}
\end{question}
\answer{4}
